% ============ 11. Ética y protección de datos ============
\section{Ética y protección de datos personales en registros de ciencia}

\subsection{Qué dice la literatura}

El tratamiento de registros administrativos de investigadores implica datos
personales y, con frecuencia, \textbf{categorías sensibles} (etnia, condición
de discapacidad, condición de víctima del conflicto). La literatura sobre
privacidad de datos muestra que la \textbf{anonimización ingenua no es
suficiente}: los ataques de reidentificación tienen éxito incluso sobre datos
incompletos o con seudónimos \citep{rocher2019estimating}, y el riesgo crece
cuando se combinan múltiples variables cuasi-identificadoras (departamento,
área, edad, cargo). Las prácticas recomendadas incluyen: minimización
(recolectar solo lo necesario), seudonimización y anonimización robusta
(k-anonimato y variantes), agregación de microdatos antes de publicar, y
análisis de riesgo de reidentificación previo a cualquier liberación de datos.
En Colombia, el marco normativo es la \textbf{Ley 1581 de 2012} (protección de
datos personales) y sus decretos reglamentarios, con reglas específicas para
datos sensibles (autorización explícita, finalidad, seguridad).

\subsection{Relación con este repositorio}

El proyecto procesa el padrón de MinCiencias con variables sensibles
(auto-declaradas): grupo étnico, discapacidad y condición de víctima, además de
identificadores de persona (\texttt{id\_persona\_pr}) y ubicación. Buena
práctica ya adoptada: la construcción de una "gran tabla sin identificadores"
para análisis y la publicación de agregados (evidencias y figuras). La
auditoría identificó riesgos a formalizar: (a) el consolidado versionado
(XLSX/CSV) incluye \texttt{ID\_PERSONA\_PR}, edad y geolocalización — microdatos
personales en un repositorio público; (b) no hay documento de tratamiento de
datos personales ni análisis formal de riesgo de reidentificación de la "gran
tabla"; (c) la descarga pública por convocatoria de la fuente (Socrata) es
legítima, pero el proyecto debe declarar su régimen de tratamiento y retención.

\subsection{Mejoras recomendadas}

\begin{enumerate}
  \item Publicar una \textbf{nota de tratamiento de datos personales} (Ley
        1581/2012): finalidad, base de legitimación, categorías sensibles,
        medidas de seguridad y política de retención del consolidado.
  \item Realizar un \textbf{análisis de riesgo de reidentificación} de la
        "gran tabla" y de los microdatos publicados (estilo
        \citet{rocher2019estimating}); si el riesgo es no despreciable,
        agregar o generalizar las variables cuasi-identificadoras.
  \item Mantener la política de \textbf{no publicar microdatos identificables}
        (ya correcta en el diseño) y documentarla explícitamente en el README y
        en el catálogo, indicando qué archivos son de acceso restringido.
\end{enumerate}
